\section{实验结果}\label{sec:task2-results}

本节是实际运行结果的登记位置。当前文本不把待采集的照片、待核对的标定板尺寸或未运行程序得到的数字视为实验结果；完成采集和运行后，应直接依据程序输出和保存文件替换表\ref{tab:task2-results}中的占位。

\begin{table}[htbp]
    \centering
    \caption{手机相机标定结果}
    \label{tab:task2-results}
    \begin{tabular}{lll}
        \hline
        类别 & 指标 & 实际值 \\
        \hline
        数据 & 原始图像总数 & 待实际运行填入 \\
        数据 & 有效标定图像数 & 待实际运行填入 \\
        数据 & 角点检测失败数 & 待实际运行填入 \\
        数据 & 图像尺寸（宽 $\times$ 高） & 待实际运行填入 \\
        内参 & 相机矩阵 \(\mathbf K\) & 待实际运行填入 \\
        内参 & \(f_x\) & 待实际运行填入 \\
        内参 & \(f_y\) & 待实际运行填入 \\
        内参 & \(c_x\) & 待实际运行填入 \\
        内参 & \(c_y\) & 待实际运行填入 \\
        畸变 & \((k_1,k_2,p_1,p_2,k_3,\ldots)\) & 待实际运行填入 \\
        误差 & OpenCV RMS（像素） & 待实际运行填入 \\
        误差 & 平均重投影误差（像素） & 待实际运行填入 \\
        \hline
    \end{tabular}
\end{table}

相机矩阵应按式\ref{eq:task2-intrinsic}的三阶矩阵形式抄录，畸变系数须注明 OpenCV 返回的完整长度和顺序。平均重投影误差的计算口径应与式\ref{eq:task2-reprojection-error}一致。若输出中包含每幅图像的误差，还应在附表中保留文件名、角点检测状态和单幅误差，便于定位异常照片。

\subsection{结果图与文件记录}

若程序生成了角点绘制图或畸变校正图，应在最终集成时插入本节，并在正文中说明其输入文件和生成方式。没有实际生成的图像不得用示意图或网络图片代替实验结果。参数文件、运行日志和图像清单应与表中数值保持一一对应。
